\documentclass[11pt,a4paper]{article}
\usepackage[utf8]{inputenc}
\usepackage[T1]{fontenc}
\usepackage[french]{babel}
\usepackage[top=2cm,bottom=2cm,left=2cm,right=2cm]{geometry}
\usepackage{xcolor}
\usepackage{tcolorbox}
\tcbuselibrary{skins,breakable}
\usepackage{fancyhdr}
\usepackage{enumitem}
\usepackage{lmodern}

% --- Couleur discipline : ÉCONOMIE ---
\definecolor{mainColor}{RGB}{0,80,160}
\definecolor{darkGray}{RGB}{50,50,50}

\tcbset{boxrule=0.5pt, arc=3pt, left=5pt, right=5pt, top=3pt, bottom=3pt, breakable}

\newtcolorbox{defbox}[1]{
  colback=mainColor!8, colframe=mainColor, coltitle=white,
  fonttitle=\bfseries\footnotesize, title=DÉFINITION \textemdash\ #1}

\newtcolorbox{keybox}{
  colback=darkGray!7, colframe=darkGray!50, coltitle=white,
  fonttitle=\bfseries\footnotesize, title=POINT CLÉ}

\newtcolorbox{exbox}{
  colback=mainColor!5, colframe=mainColor!40,
  fonttitle=\bfseries\footnotesize, title=EXEMPLE}

\pagestyle{fancy}\fancyhf{}
\fancyhead[L]{\small\textcolor{mainColor}{\textbf{CEJM — BTS SIO}}}
\fancyhead[C]{\small\textbf{Chapitre 10 — Les choix de production}}
\fancyhead[R]{\small\textcolor{mainColor}{\textbf{ÉCONOMIE}}}
\fancyfoot[C]{\small\thepage}
\renewcommand{\headrulewidth}{0.5pt}
\setlength{\headheight}{15pt}
\setlist{itemsep=2pt, topsep=3pt, parsep=0pt}

\begin{document}

\begin{center}
  {\Large\bfseries\textcolor{mainColor}{Chapitre 10}}\\[2pt]
  {\large Les choix de production de l'entreprise}\\[4pt]
  \textit{\textcolor{darkGray}{Comment l'entreprise décide-t-elle de produire et à quel coût ?}}
\end{center}
\vspace{4pt}\hrule\vspace{8pt}

\section*{Définitions essentielles}

\begin{defbox}{Facteurs de production}
Ressources utilisées par l'entreprise pour produire : le \textbf{travail} (RH), le \textbf{capital} (machines, bâtiments) et les \textbf{ressources naturelles} (matières premières, énergie).
\end{defbox}

\vspace{4pt}
\begin{defbox}{Coûts fixes (CF)}
Charges qui \textbf{ne varient pas} avec le niveau de production (loyer, amortissements, salaires de direction).
\end{defbox}

\vspace{4pt}
\begin{defbox}{Coûts variables (CV)}
Charges qui \textbf{varient proportionnellement} avec le niveau de production (matières premières, énergie, main-d'œuvre de production).
\end{defbox}

\vspace{4pt}
\begin{defbox}{Seuil de rentabilité (point mort)}
Niveau de \textbf{chiffre d'affaires} à partir duquel l'entreprise commence à réaliser un bénéfice (CA qui couvre tous les coûts).
\end{defbox}

\vspace{4pt}
\begin{defbox}{Productivité}
Rapport entre la \textbf{production obtenue} et les \textbf{facteurs utilisés}. Elle mesure l'efficacité de la combinaison productive.
\end{defbox}

\section*{Les facteurs de production}

\begin{center}
\begin{tabular}{|l|l|l|}
\hline
\textbf{Facteur} & \textbf{Contenu} & \textbf{Rémunération} \\
\hline
\textbf{Travail} & Ressources humaines, compétences & Salaires \\
\textbf{Capital technique} & Machines, outillage, bâtiments & Amortissement \\
\textbf{Capital financier} & Fonds investis & Intérêts / dividendes \\
\textbf{Ressources naturelles} & Matières premières, énergie & Prix de marché \\
\hline
\end{tabular}
\end{center}

\begin{keybox}
L'entreprise cherche la \textbf{combinaison productive optimale} : la quantité de travail et de capital qui minimise le coût pour une production donnée.
\end{keybox}

\section*{La structure des coûts}

\subsection*{Formules essentielles}
\begin{itemize}
  \item \textbf{Coût total (CT)} = Coûts fixes (CF) + Coûts variables (CV)
  \item \textbf{Coût moyen (CU)} = CT / Quantité produite
  \item \textbf{Marge sur coût variable (MCV)} = CA $-$ CV = CF + Résultat
  \item \textbf{Taux de MCV} = MCV / CA × 100
  \item \textbf{Seuil de rentabilité} = CF / Taux de MCV
\end{itemize}

\begin{exbox}
Une entreprise a CF = 50 000 €, CV unitaire = 10 €, Prix = 25 €. \\
MCV unitaire = 25 $-$ 10 = 15 €. Seuil de rentabilité = 50 000 / 15 = \textbf{3 333 unités à vendre}.
\end{exbox}

\section*{La combinaison productive : Faire ou Faire faire ?}

\begin{center}
\begin{tabular}{|l|l|l|}
\hline
 & \textbf{Internalisation (Faire)} & \textbf{Externalisation (Faire faire)} \\
\hline
\textbf{Avantages} & Contrôle qualité, savoir-faire & Réduction des coûts fixes, flexibilité \\
\textbf{Inconvénients} & Coûts fixes élevés, rigidité & Dépendance fournisseur, risque qualité \\
\textbf{Exemples} & Production en propre & Sous-traitance, cloud computing \\
\hline
\end{tabular}
\end{center}

\begin{keybox}
L'entreprise externalise les activités \textbf{non stratégiques} pour se concentrer sur son \textbf{cœur de métier}.
\end{keybox}

\section*{La productivité}

\begin{itemize}
  \item \textbf{Productivité du travail} = Production / Nombre d'heures travaillées (ou nb de salariés)
  \item \textbf{Productivité du capital} = Production / Capital utilisé
  \item \textbf{Gains de productivité} → $\downarrow$ coûts unitaires → $\uparrow$ compétitivité → $\uparrow$ profit ou $\downarrow$ prix
\end{itemize}

\begin{keybox}
Le \textbf{progrès technique} (automatisation, IA, organisation) est le principal moteur des gains de productivité.
\end{keybox}

\section*{Les contraintes de production}

\begin{itemize}
  \item \textbf{Coût des facteurs} : hausse du prix de l'énergie, des matières premières ou des salaires → $\uparrow$ coûts variables
  \item \textbf{Financement} : capacité à investir dans le capital technique
  \item \textbf{Réglementation} : normes environnementales, droit du travail
  \item \textbf{Conjoncture} : demande incertaine → décision de produire ou d'arrêter
\end{itemize}

\begin{exbox}
\textbf{Duralex} (2022) : hausse du coût de l'énergie de +50\,\% → arrêt de la production pendant 5 mois. Les salariés ont été placés en chômage partiel.\\
\textbf{Danone} : investissement de 43 M€ dans son usine du Gers pour reconvertir la production vers les protéines végétales.
\end{exbox}

\section*{Points clés à retenir}

\begin{keybox}
L'entreprise cherche la \textbf{combinaison productive optimale} : produire au moindre coût avec les facteurs disponibles.
\end{keybox}
\vspace{3pt}
\begin{keybox}
Le \textbf{seuil de rentabilité} est un indicateur crucial : en dessous, l'entreprise perd de l'argent.
\end{keybox}
\vspace{3pt}
\begin{keybox}
\textbf{Faire vs faire faire} : l'externalisation réduit les coûts fixes mais crée une dépendance fournisseur.
\end{keybox}
\vspace{3pt}
\begin{keybox}
La \textbf{hausse des coûts} (énergie, matières premières) peut remettre en cause la viabilité économique de la production.
\end{keybox}

\section*{Mots-clés}
\textcolor{mainColor}{\textbf{Facteurs de production}} $\cdot$
\textcolor{mainColor}{\textbf{Coûts fixes / variables}} $\cdot$
\textcolor{mainColor}{\textbf{Coût total / moyen}} $\cdot$
\textcolor{mainColor}{\textbf{Marge sur coût variable}} $\cdot$
\textcolor{mainColor}{\textbf{Seuil de rentabilité}} $\cdot$
\textcolor{mainColor}{\textbf{Productivité}} $\cdot$
\textcolor{mainColor}{\textbf{Internalisation / Externalisation}} $\cdot$
\textcolor{mainColor}{\textbf{Combinaison productive}} $\cdot$
\textcolor{mainColor}{\textbf{Progrès technique}}

\end{document}
